The deterministic method uses a changing regularization parameter to make
each remaining correction diffuse. It then accelerates that correction on
an explicit convex set. The set contains the unknown correction optimum,
but its definition requires only the current source and the degrees.

\subsection{A diffuse-source stage}

Write
\begin{equation}
 \vomega:=\mD^{1/2}\one,
 \qquad \lambda_r:=\alpha r.
 \label{eq:det-degree-weight}
\end{equation}
Here $r\geq\rho$ is the stage regularizer; $\rho$ remains the final target.
At the start of a stage we have a sparse baseline $\overline\vx$ satisfying
\begin{equation}
 \vzero\leq\overline\vx\leq\vx_r^*,\qquad
 \vh:=\vb-\mQ\overline\vx,\qquad
 \vzero\leq\vh\leq4\lambda_r\vomega.
 \label{eq:det-stage-invariant}
\end{equation}
The vector $\vh$ is a residual source, distinct from the seed distribution
$\vs$. Its weighted mass is
\begin{equation}
 m_h:=\vomega^\trans\vh
      =\alpha(1-\vomega^\trans\overline\vx),\qquad
 m_r:=m_h/\alpha\leq1.
 \label{eq:det-source-mass}
\end{equation}
We use two comparison vectors only in the analysis:
\begin{equation}
 \vxi^*:=\vx_r^*-\overline\vx,
 \qquad \vt:=\mQ^{-1}\vh=\vx_0^*-\overline\vx.
 \label{eq:det-comparators}
\end{equation}
The algorithm computes neither inverse nor optimum.

\begin{lemma}[An explicit correction domain]
\label{lem:det-correction-domain}
Both $\vxi^*$ and $\vt$ belong to
\begin{equation}
 \gK_r:=\{\vu:\vzero\leq\vu\leq4r\vomega,
                    \ \vomega^\trans\vu\leq m_r\}.
 \label{eq:det-correction-domain}
\end{equation}
In fact $\vzero\leq\vxi^*\leq\vt\leq4r\vomega$ and
$\vomega^\trans\vt=m_r$.
\end{lemma}
\begin{proof}
The order follows from \cref{eq:rppr-bias}, the baseline invariant,
inverse positivity, and
$\mQ^{-1}\vomega=\vomega/\alpha$:
$\vt\leq4\alpha r\mQ^{-1}\vomega=4r\vomega$.
Multiplying $\mQ\vt=\vh$ by $\vomega^\trans$ proves the mass identity.
\end{proof}

The correction objective is the smooth quadratic
\begin{equation}
 J_r(\vxi):=\tfrac12\vxi^\trans\mQ\vxi
                 -(\vh-\lambda_r\vomega)^\trans\vxi.
 \label{eq:det-correction-objective}
\end{equation}
On nonnegative corrections it satisfies
$J_r(\vxi)=F_r(\overline\vx+\vxi)-F_r(\overline\vx)$.
Thus $\vxi^*$ minimizes $J_r$ on $\gK_r$, and correction gaps are
exactly the desired gaps for the full candidate.

Choose a dyadic $\theta$ by halving $1/2$ until $\theta^2\leq\alpha$, and set
\begin{equation}
 \mu_{\rm c}:=\theta^2,\qquad \chi:=1-\theta,
 \qquad \alpha/4\leq\mu_{\rm c}\leq\alpha.
 \label{eq:det-acceleration-parameters}
\end{equation}
The subscript distinguishes $\mu_{\rm c}$ from a source paper's curvature
notation. Start with $\vxi_0=\vz_0=\vzero$ and iterate
\begin{equation}
 \begin{aligned}
 \vy_k&=\frac{\vxi_k+\theta\vz_k}{1+\theta},\\
 \vq_k&=\chi\vz_k+\theta\vy_k
       -\frac{\mQ\vy_k-\vh+\lambda_r\vomega}{\theta},\\
 \vz_{k+1}&=\operatorname{Proj}_{\gK_r}(\vq_k),\\
 \vxi_{k+1}&=\chi\vxi_k+\theta\vz_{k+1}.
 \end{aligned}
 \label{eq:det-iteration}
\end{equation}
Projection is Euclidean. We call $\vz_k$ the \emph{kinetic vector} to
distinguish the coordinates updated by the projection from the longer-lived
primal correction $\vxi_k$. Both stay in $\gK_r$ by convexity.

\paragraph{Why both constraints are present.}
The upper box controls the sign of a projection normal at a saturated
coordinate. The mass cap controls the common normal introduced by the
weighted sum constraint. Together they allow the \emph{same} Euclidean
projection to satisfy a comparison inequality in the $\mQ$ metric
(\cref{lem:det-sector}). That second comparison is the additional ingredient
needed to bound local work. Neither constraint requires knowing the support
of $\vx_r^*$, and the kinetic support need not be contained in it.

\paragraph{Explicit projection.}
There is a scalar $\gamma\geq0$ such that
\begin{equation}
 \frac{z_{k+1,i}}{\omega_i}
 =\min\{4r,(q_{k,i}/\omega_i-\gamma)_+\}.
 \label{eq:det-waterfill}
\end{equation}
If the right side at $\gamma=0$ has weighted mass at most $m_r$, use zero;
otherwise choose $\gamma$ so that $\vomega^\trans\vz_{k+1}=m_r$.
These are exactly the projection KKT conditions. The finite threshold search
and the cost of reporting its positive coordinates are proved in
\cref{sec:deterministic-implementation}.

\subsection{A certified stage handoff}

A small objective gap does not by itself make the candidate a subsolution.
We use one projected-gradient step and a downward density correction at the
\emph{end} of a stage. Intermediate iterations remain precisely
\cref{eq:det-iteration}.

\begin{lemma}[Safe repair]
\label{lem:det-repair}
Let $0<\delta\leq\lambda_r/2$. Suppose a nonnegative candidate
$\widetilde\vx$ has $F_r(\widetilde\vx)-F_r(\vx_r^*)\leq\tau$, where
$\tau\leq\alpha\delta^2/8$. Let the old baseline satisfy
$\vzero\leq\overline\vx\leq\vx_r^*$ and
$\mQ\overline\vx\leq\vb$; it may be zero. Define
\begin{equation}
 \vp=[\widetilde\vx-(\mQ\widetilde\vx-\vb+\lambda_r\vomega)]_+,
 \qquad \vu=[\vp-\delta\vomega]_+.
 \label{eq:det-repair-map}
\end{equation}
Then $\overline\vx^+:=\max\{\overline\vx,\vu\}$ satisfies
\begin{equation}
 \begin{gathered}
 \vzero\leq\overline\vx\leq\overline\vx^+\leq\vx_r^*,\qquad
 \vzero\leq\vb-\mQ\overline\vx^+\leq2\lambda_r\vomega,\\
 \vol(\supp\overline\vx^+)\leq1/r,\qquad
 F_r(\overline\vx^+)-F_r(\vx_r^*)\leq2\delta^2/r.
 \end{gathered}
 \label{eq:det-repaired-invariant}
\end{equation}
The coordinatewise maximum can be omitted if monotone baselines are not
needed.
\end{lemma}
\begin{proof}
Write $\vdelta=\widetilde\vx-\vp$. Projection optimality and
$\mQ\preceq\mI$ give
\begin{equation}
 F_r(\widetilde\vx)-F_r(\vp)\geq\tfrac12\norm{\vdelta}_2^2,
 \qquad (\mI-\mQ)\vdelta\in
 \partial(\phi_r+I_{\R_+^n})(\vp).
 \label{eq:det-pg-certificate}
\end{equation}
For the second assertion, the projection normal is
$\vdelta-\nabla\phi_r(\widetilde\vx)$; adding
$\nabla\phi_r(\vp)$ gives $(\mI-\mQ)\vdelta$.
Strong convexity and descent imply
\[
 \norm{\vp-\vx_r^*}_2\leq\sqrt{2\tau/\alpha}\leq\delta/2,
 \qquad \norm{(\mI-\mQ)\vdelta}_2\leq\sqrt{2\tau}\leq\delta/2.
\]
At every positive coordinate of $\vp$ the normal cone vanishes. Hence
$(\mQ\vp-\vb)_i\leq\delta/2-\lambda_r\omega_i\leq0$.
At a zero coordinate the same inequality follows from the Stieltjes signs.
Clipping gives $\vzero\leq\vu\leq\vx_r^*$ and
$\vzero\leq\vx_r^*-\vu\leq2\delta\vomega$, using $\omega_i\geq1$.
If $u_i>0$, then $u_i=p_i-\delta\omega_i$ and
$\vu\geq\vp-\delta\vomega$. Therefore
$(\mQ\vu)_i\leq(\mQ\vp)_i-\alpha\delta\omega_i\leq b_i$;
zero coordinates again follow by the off-diagonal sign.

The maximum of two nonnegative subsolutions is a subsolution: on row $i$,
choose the vector that attains the maximum there; the diagonal term is
unchanged and every other coordinate can only increase. Both vectors are
below $\vx_r^*$, and the maximum retains the displayed error bound.
The identity $|\mQ|\vomega=\vomega$ then gives
\[
 \vb-\mQ\overline\vx^+
 \leq\vb-\mQ\vx_r^*+2\delta\vomega
 \leq(\lambda_r+2\delta)\vomega\leq2\lambda_r\vomega.
\]
Support containment gives the volume bound and makes the linear KKT term
in the objective expansion zero. Thus the gap is
$\tfrac12\norm{\overline\vx^+-\vx_r^*}_{\mQ}^2
 \leq\tfrac12(2\delta)^2\vol(\mathcal S_r^*)\leq2\delta^2/r$.
\end{proof}

\subsection{The complete deterministic algorithm}

\begin{algorithm}[t]
\caption{Deterministic accelerated local RPPR}
\label{alg:deterministic-rppr}
\begin{algorithmic}[1]
\REQUIRE Degree and adjacency-list access, seed $v$, $\alpha,\rho,\epsobj>0$
\STATE Query $d_v$. If $\rho d_v\geq1$ or $\epsobj\geq\alpha/2$, return $\vzero$.
\STATE If $\alpha=1$, return the single coordinate
       $\widehat x_v=(1-\rho d_v)/\sqrt{d_v}$.
\STATE Choose $\theta,\mu_{\rm c},\chi$ as in
       \cref{eq:det-acceleration-parameters}; set
       $r_{\rm old}=1/d_v$ and $\overline\vx=\vzero$.
\REPEAT
  \STATE $r\gets\max\{\rho,r_{\rm old}/2\}$; set $\delta\gets\alpha r/2$.
  \IF{$r=\rho$}
    \STATE Halve $\delta$ until $2\delta^2/r\leq\epsobj$.
  \ENDIF
  \STATE Assemble $\vh=\vb-\mQ\overline\vx$ locally; compute $m_r$.
  \STATE Set $\vxi=\vz=\vzero$, $a_0=1$, and $\tau=\alpha\delta^2/8$.
  \WHILE{$a_0>\tau$}
    \STATE Apply \cref{eq:det-iteration} using the sparse threshold reporter.
    \STATE $a_0\gets\chi a_0$.
  \ENDWHILE
  \STATE Materialize $\widetilde\vx=\overline\vx+\vxi$; apply
         \cref{eq:det-repair-map} and set
         $\overline\vx\gets\max\{\overline\vx,\vu\}$.
  \STATE $r_{\rm old}\gets r$.
\UNTIL{$r=\rho$}
\RETURN The sparse coordinate list of $\overline\vx$.
\end{algorithmic}
\end{algorithm}

The scalar $a_0$ is a computable upper bound on the current comparison
energy, as proved next. In particular, the stopping test uses no optimum,
global gradient, or unknown complementarity margin.
The first stage has $r\geq1/(2d_v)$, so its source obeys
$\vb\leq2\alpha r\vomega$. At subsequent stages,
\cref{eq:det-repaired-invariant} and $r\geq r_{\rm old}/2$ imply
$\vh\leq4\alpha r\vomega$. Monotonicity of $\vx_r^*$ as $r$ decreases
preserves baseline safety. These facts close the continuation induction.
